\begin{ouverture}[geometrie]{14}{Construire et comprendre les triangles}{Partie 3 · Espace et géométrie}
{Angles · Médiatrices · Cercle circonscrit · Hauteurs · Médianes}
\objectif{Construire un triangle à partir de données.}
\objectif{Tracer hauteurs, médianes et médiatrices.}
\objectif{Calculer des aires et expliquer une preuve.}
\prerequis{Angles, compas, règle et équerre.}
\end{ouverture}
\pagestyle{otMain}\essentiel\fontsize{11}{15}\selectfont
\phantomsection\label{cours-14-1}
\essSection{1}{Angles et construction}
\begin{propriete}La somme des angles d’un triangle est $180^\circ$. Pour le voir, tracer par un sommet la parallèle au côté opposé : les angles alternes internes reportent les deux autres angles à ce sommet, où ils forment ensemble un angle plat.\end{propriete}
\label{fig-cours-14-1-cours-triangle-angles}\figureGeo{cours-triangle-angles}

\begin{savoirfaire}[Choisir les instruments]On peut construire un triangle avec trois longueurs, deux côtés et l’angle compris, ou un côté et deux angles. Avec trois côtés, tracer d’abord un segment puis deux arcs de compas. Pour un triangle non aplati, le plus grand côté doit être strictement inférieur à la somme des deux autres.\end{savoirfaire}
\begin{exemple}4, 5 et 7 cm permettent un triangle car $7<4+5$. 2, 3 et 6 cm ne le permettent pas car $6>2+3$.\end{exemple}
\clearpage
\phantomsection\label{cours-14-2}
\essSection{2}{Médiatrices et cercle circonscrit}
\begin{definition}La médiatrice d’un segment est la droite perpendiculaire à ce segment passant par son milieu. Ses points sont à égale distance des deux extrémités, et tout point équidistant des extrémités appartient à la médiatrice.\end{definition}
\label{fig-cours-14-2-cours-cercle-circonscrit}\figureGeo{cours-cercle-circonscrit}

\begin{propriete}Les trois médiatrices d’un triangle sont concourantes. Leur point commun O est équidistant des trois sommets ; il est le centre du cercle circonscrit. Pour construire O, deux médiatrices suffisent ; on contrôle avec le troisième sommet.\end{propriete}
\begin{savoirfaire}[Triangles particuliers]Dans un triangle isocèle, la droite du sommet principal au milieu de la base est un axe de symétrie : elle porte la médiatrice de la base, la hauteur et la médiane correspondantes. Dans un équilatéral, cette propriété vaut aux trois sommets.\end{savoirfaire}
\begin{propriete}Dans un triangle rectangle, on observe que le centre du cercle circonscrit est le milieu de l’hypoténuse, le côté opposé à l’angle droit : le triangle est inscrit dans un demi-cercle dont l’hypoténuse est un diamètre. Cette propriété est admise en Cinquième ; elle sera démontrée en Quatrième.\end{propriete}
\clearpage
\phantomsection\label{cours-14-3}
\essSection{3}{Tracer les hauteurs}
\begin{definition}Une hauteur est une droite passant par un sommet et perpendiculaire à la droite qui porte le côté opposé. Pour calculer une aire, on utilise aussi la longueur du segment perpendiculaire entre le sommet et cette droite.\end{definition}
\label{fig-cours-14-3-cours-hauteurs}\figureGeo{cours-hauteurs}

\begin{propriete}Les trois hauteurs sont concourantes en un point appelé orthocentre ; cette propriété, constatée sur des figures, est admise. Il est à l’intérieur pour un triangle aigu, au sommet de l’angle droit pour un triangle rectangle et à l’extérieur pour un triangle obtus.\end{propriete}
\begin{attention}Une hauteur peut rencontrer le prolongement d’un côté. La médiane et la médiatrice ne sont pas définies par la même condition : la hauteur exige la perpendicularité, pas le milieu.\end{attention}
\clearpage
\phantomsection\label{cours-14-4}
\essSection{4}{Aires et médianes}
\begin{propriete}L’aire d’un triangle est $A=\frac{b\times h}{2}$, où $h$ est la hauteur relative à la base $b$. Convertir les longueurs dans la même unité avant de calculer.\end{propriete}
\label{fig-cours-14-4-cours-mediane}\figureGeo{cours-mediane}

\begin{definition}Une médiane est la droite passant par un sommet et par le milieu du côté opposé. Le segment joignant ces deux points partage le triangle en deux triangles d’aires égales : leurs bases ont la même longueur et ils ont la même hauteur. On observe que les trois médianes sont concourantes, en un point appelé centre de gravité ; cette propriété est admise.\end{definition}
\begin{exemple}Pour une base de 8 cm et une hauteur de 3 cm, l’aire vaut $8\times3:2=12$ cm$^2$. La médiane issue du sommet correspondant donne deux aires de 6 cm$^2$.\end{exemple}
\begin{attention}Dans un triangle quelconque, une médiane n’est pas nécessairement perpendiculaire au côté. Un schéma codé donne les hypothèses ; l’apparence ne constitue pas une preuve.\end{attention}

\clearpage\rubrique{Découverte}
\begin{blocLivre}{ex-14-1}\begin{exercice}[1]{Troisième angle}
Deux angles d’un triangle mesurent $48^\circ$ et $67^\circ$. Calculer le troisième.
\label{fig-ex-14-1}\figureGeo{tri-1}
\end{exercice}\end{blocLivre}
\begin{blocLivre}{ex-14-2}\begin{exercice}[1]{Un isocèle}
ABC est isocèle en A et $\widehat A=44^\circ$. Calculer les angles en B et C.
\label{fig-ex-14-2}\figureGeo{tri-2}
\end{exercice}\end{blocLivre}
\begin{blocLivre}{ex-14-3}\begin{exercice}[1]{Une construction à trois côtés}
Construire ABC avec $AB=6$ cm, $AC=4$ cm et $BC=5$ cm.
\label{fig-ex-14-3}\figureGeo{tri-3}
\end{exercice}\end{blocLivre}
\begin{blocLivre}{ex-14-4}\begin{exercice}[1]{Aire directe}
Un triangle a une base de 9 cm et une hauteur correspondante de 4 cm. Calculer son aire.
\label{fig-ex-14-4}\figureGeo{tri-4}
\end{exercice}\end{blocLivre}
\begin{blocLivre}{ex-14-5}\begin{exercice}[1]{Reconnaître une médiane}
M est le milieu de [BC]. Pourquoi (AM) est-elle une médiane du triangle ABC ?
\label{fig-ex-14-5}\figureGeo{tri-5}
\end{exercice}\end{blocLivre}

\rubrique{Entraînement}
\begin{blocLivre}{ex-14-6}\begin{exercice}[2]{Deux côtés et un angle}
Construire $AB=5$ cm, $AC=3$ cm et $\widehat{BAC}=60^\circ$.
\label{fig-ex-14-6}\figureGeo{tri-6}
\end{exercice}\end{blocLivre}
\begin{blocLivre}{ex-14-7}\begin{exercice}[2]{Un côté et deux angles}
Construire $AB=6$ cm, $\widehat{BAC}=50^\circ$ et $\widehat{ABC}=60^\circ$. Calculer $\widehat C$.
\label{fig-ex-14-7}\figureGeo{tri-7}
\end{exercice}\end{blocLivre}
\begin{blocLivre}{ex-14-8}\begin{exercice}[2]{Hauteur extérieure}
Dans le triangle obtus donné, tracer la hauteur issue de C relative à (AB). Pourquoi faut-il prolonger [AB] ?
\label{fig-ex-14-8}\figureGeo{tri-8}
\end{exercice}\end{blocLivre}
\begin{blocLivre}{ex-14-9}\begin{exercice}[2]{Cercle circonscrit}
Construire le cercle passant par les trois sommets du triangle donné.
\label{fig-ex-14-9}\figureGeo{tri-9}
\end{exercice}\end{blocLivre}
\begin{blocLivre}{ex-14-10}\begin{exercice}[2]{Hauteur et aire}
Un triangle a une aire de 24 cm$^2$ et une base de 8 cm. Trouver la hauteur correspondante.
\label{fig-ex-14-10}\figureGeo{tri-10}
\end{exercice}\end{blocLivre}

\rubrique{Problèmes}
\begin{blocLivre}{ex-14-11}\begin{exercice}[3]{Aire et médiane}
ABC a une aire de 35 cm$^2$ et M est le milieu de [BC]. Calculer les aires de ABM et ACM en justifiant.
\label{fig-ex-14-11}\figureGeo{tri-11}
\end{exercice}\end{blocLivre}
\begin{blocLivre}{ex-14-12}\begin{exercice}[3]{Triangle rectangle : hauteurs et cercle}
ABC est rectangle en A. Identifier deux de ses hauteurs et le point de concours des trois hauteurs. Construire ensuite le cercle circonscrit à ABC : où se trouve son centre ?
\label{fig-ex-14-12}\figureGeo{tri-12}
\end{exercice}\end{blocLivre}
\begin{blocLivre}{ex-14-13}\begin{exercice}[3]{Un isocèle et ses droites}
ABC est isocèle en A et M est le milieu de [BC]. Expliquer pourquoi (AM) est à la fois médiane, médiatrice de [BC] et hauteur. Pourquoi le centre du cercle circonscrit appartient-il à (AM) ?
\label{fig-ex-14-13}\figureGeo{tri-13}
\end{exercice}\end{blocLivre}
\begin{blocLivre}{ex-14-14}\begin{exercice}[3]{Un tracé impossible}
Peut-on construire un triangle de côtés 3 cm, 4 cm et 8 cm ? Et de côtés 3, 4 et 7 cm ?
\label{fig-ex-14-14}\figureGeo{tri-14}
\end{exercice}\end{blocLivre}
\begin{blocLivre}{ex-14-15}\begin{exercice}[3]{Deux bases pour la même aire}
ABC a $AB=12$ cm, la hauteur issue de C vaut 5 cm et $BC=10$ cm. Trouver la hauteur issue de A puis l’aire de chaque triangle créé par la médiane issue de C.
\label{fig-ex-14-15}\figureGeo{tri-15}
\end{exercice}\end{blocLivre}

\clearpage
